% !PRINTED 83
\weberpage{83}
%<<CONT
Wäre also $J$ nicht durch $2^{\lambda-2}$ theilbar, so müssten alle $\beta_2$ durch 2 theilbar sein, was wieder gegen die Forderung IV ist.

e) \gesperrt{Wenn $q$ eine einfach in $m$ aufgehende Primzahl ist, so muss $q-1$ wenigstens durch eine der in $J$ aufgehenden Primzahlen theilbar sein.}

Denn wäre $c_1=q-1$ relativ prim zu $J$, so müssten alle $\beta_1$ durch $q-1$ theilbar sein, im Widerspruch mit IV. Die weiteren Bedingungen für $m$ ergeben sich später.

\gesperrt{Bestimmung der Gruppe $\mathfrak{B}$:}

Die Gruppe $\mathfrak{B}$ ist bestimmt, wenn ihre Basis $g_1, g_2, \dots g_\nu$ gegeben ist. Die Elemente der Basis wollen wir durch ihre Indices bestimmen, und führen also folgende Bezeichnung ein. Es seien
\weq{7}{
\begin{aligned}
&\gamma_{1,1}, \gamma_{1,2}, \dots \gamma_{1,\mu} \quad \text{die Indices von } g_1\\
&\gamma_{2,1}, \gamma_{2,2}, \dots \gamma_{2,\mu} \quad \text{\textquotedbl{}} \quad \text{\textquotedbl{}} \quad \text{\textquotedbl{}} \quad g_2\\
&\dots\dots\dots\dots\dots\dots\dots\dots\dots\dots\dots\dots\dots\dots\\
&\gamma_{\nu,1}, \gamma_{\nu,2}, \dots \gamma_{\nu,\mu} \quad \text{\textquotedbl{}} \quad \text{\textquotedbl{}} \quad \text{\textquotedbl{}} \quad g_\nu.
\end{aligned}
}
Die Indices $\beta_1, \beta_2, \dots \beta_\mu$ eines beliebigen Elementes $b$ erhält man dann nach (3) in der Form:
\weq{8}{
\begin{aligned}
\beta_1 &\equiv \gamma_{1,1}x_1+\gamma_{2,1}x_2+\dots+\gamma_{\nu,1}x_\nu\ (\mathrm{mod}\ c_1)\\
\beta_2 &\equiv \gamma_{1,2}x_1+\gamma_{2,2}x_2+\dots+\gamma_{\nu,2}x_\nu\ (\mathrm{mod}\ c_2)\\
&\dots\dots\dots\dots\dots\dots\dots\dots\dots\dots\dots\dots\dots\dots\\
\beta_\mu &\equiv \gamma_{1,\mu}x_1+\gamma_{2,\mu}x_2+\dots+\gamma_{\nu,\mu}x_\nu\ (\mathrm{mod}\ c_\mu),
\end{aligned}
}
und wir fragen nun nach den nothwendigen und hinreichenden Bedingungen für die Zahlen $\gamma_{h,k}$, damit $g_1, g_2 \dots g_\nu$ die Basis einer Gruppe $\mathfrak{B}$ von den verlangten Eigenschaften ist. Nach 1. ist hierfür zuerst nothwendig:

\gesperrt{2. Die Zahlen $\gamma_{h,k}$ müssen so beschaffen sein, dass die Congruenzen}
\weq{9}{
\begin{aligned}
\gamma_{1,1}x_1+\gamma_{2,1}x_2+\dots+\gamma_{\nu,1}x_\nu &\equiv 0\ (\mathrm{mod}\ c_1)\\
\gamma_{1,2}x_1+\gamma_{2,2}x_2+\dots+\gamma_{\nu,2}x_\nu &\equiv 0\ (\mathrm{mod}\ c_2)\\
&\dots\dots\dots\dots\dots\dots\dots\dots\dots\dots\dots\dots\\
\gamma_{1,\mu}x_1+\gamma_{2,\mu}x_2+\dots+\gamma_{\nu,\mu}x_\nu &\equiv 0\ (\mathrm{mod}\ c_\mu)
\end{aligned}
}
\gesperrt{dann und nur dann erfüllt sind. wenn zugleich die Congruenzen}
% [sic: Weber prints a full stop before the lower-case ``wenn''; a comma is grammatically required.]
\weq{10}{x_1\equiv 0\ (\mathrm{mod}\ i_1),\quad x_2\equiv 0\ (\mathrm{mod}\ i_2),\quad \dots\ x_\nu\equiv 0\ (\mathrm{mod}\ i_\nu)}
\gesperrt{bestehen.}
